\documentclass[10pt,a4paper]{amsart}
\usepackage[margin=19mm]{geometry}
\usepackage{amsmath,amssymb,mathtools}
\usepackage[scr=rsfs,cal=euler]{mathalfa}
\usepackage{xfrac}
\usepackage[unicode,hidelinks,bookmarks=false]{hyperref}
\newcommand{\ie}{i.e.\ }
\newcommand{\cf}{cf.\ }
\newcommand{\resp}{respectively\ }
\newcommand{\Subsection}{\subsection}
\providecommand{\nobreakdash}{\nobreak\hbox{-}}
\setlength{\emergencystretch}{2em}
\multlinegap0pt
\usepackage{bm}
\usepackage{bbm}
\usepackage{dsfont}
\usepackage{xspace}
\usepackage{cancel}
\usepackage{mathtools}
\usepackage{amsthm}
\makeatletter
\@ifundefined{prop}{\newtheorem{prop}{Proposition}}{}
\@ifundefined{theorem}{\newtheorem{theorem}{Theorem}}{}
\makeatother
\title[A new Representation of $\alpha$-Bernstein Operators]{A new Representation of $\alpha$-Bernstein Operators}
\author{Jamshid Saeidian, Bahareh Nouri}
\date{Source: arXiv:2407.16175v1 (CC BY 4.0)}
\begin{document}
\maketitle

\noindent\emph{Source and licence.} A new Representation of $\alpha$-Bernstein Operators. Version arXiv:2407.16175v1 (CC BY 4.0), distributed under the Creative Commons Attribution 4.0 International licence, as recorded in the arXiv licence metadata for this exact version and shown on the abstract page https://arxiv.org/abs/2407.16175v1. Full rights notice: licenses/arxiv-2407.16175-CC-BY-4.0.txt.\par\smallskip

\noindent\textit{Editorial scope.} Counted unit: the Theorem whose source label is \texttt{None} in the article (source0.tex lines 711--715, source label \texttt{None}), delivered together with the article's own complete original proof of it (source0.tex lines 717--750, one \texttt{proof} environment of 34 lines). No mathematics is written, repaired or re-derived: statement and proof are the article's own text line for line. Required context retained uncounted: rec-rel (lines 279--281); formula1 (lines 330--334); pro3.2 (lines 654--656). Every definition, display and label of those lines is retained unchanged. Omitted from this package and not redistributed: 1--278, 282--329, 335--653, 657--710, 716--716, 751--829. Nothing inside a retained excerpt is omitted or altered: every retained body is delivered exactly as the article's lines are written, so no omission receipt of any kind accompanies this package. Citations: the retained excerpts carry no citation command at all, so no bibliography entry of the article is transcribed and no reference is redistributed. Reference adaptation: none was needed and none was made; every reference inside the delivered unit resolves inside it, so no unresolved reference or citation is produced. Printed slips kept literal: no printed word, symbol, hypothesis or number of the counted statement or of its proof was altered. Layout and class: the article's own class is \texttt{elsarticle} with packages including amssymb, graphics, amsmath, amsthm, xcolor, enumerate, subfig, mathrsfs; the wrapper is the lane's pinned \texttt{amsart} editorial wrapper at 10pt on A4, which restates the theorem-like environments the retained text needs and adds the article's own macro definitions that the retained text needs (none), with extra wrapper packages (bm, bbm, dsfont, xspace, cancel, mathtools). The delivered numbering is the unit's own. Assets: no retained span contains a figure environment, a graphics-inclusion command, a PDF-page inclusion, a table or a TikZ picture; no image file is copied into this package, the package declares no asset dependency and no image file is redistributed with it. Licence: version arXiv:2407.16175v1 of this article is distributed under the Creative Commons Attribution 4.0 International licence, as recorded in the arXiv licence metadata for this exact version and shown on the abstract page https://arxiv.org/abs/2407.16175v1. A full rights notice is delivered at licenses/arxiv-2407.16175-CC-BY-4.0.txt. Characters: every retained excerpt is pure ASCII, so the delivered engine, XeLaTeX with its default Latin Modern text font, reports no missing character.
  \begin{equation}
F_{n,i}(z) = (1-z)F_{n-1,i}(z) + zF_{n-1,i-1}(z) \quad \text{for } z \in [0,1], \quad i = 0, 1, 2, \ldots, n. \label{rec-rel}
\end{equation}

\begin{eqnarray} 
F_{2,0}(z)&=&1-z+\alpha \left ( z^{2}-z \right ),\nonumber \\
F_{2,1}(z)&=&-2\alpha \left ( z^{2}-z \right ), \label{formula1}\\
F_{2,2}(z)&=&z+\alpha \left ( z^{2}-z \right ).\nonumber
\end{eqnarray} 

\begin{prop}\label{pro3.2}
The $\alpha$-Bernstein polynomials $\{F_{n,0}, F_{n,1}, \cdots, F_{n,n}\}$, defined by equations (\ref{formula1}) and (\ref{rec-rel}),  are  monotonicity preserving.
\end{prop}

\begin{theorem}
 If $q \in C[0,1] $ is a convex function, then    its $\alpha$-Bernstein operators are all convex.
 
% for $0 \leq \alpha \leq 1$, so are
\end{theorem}

\begin{proof}
We employ the method of induction to establish the result's validity, starting with the base case $n=2$.  Consider a set of real values $ \lambda_{0}, \lambda _{1} , \lambda _{2}$ forming a convex data set, ensuring $\lambda_{2} -2\lambda_{1} +\lambda_{0}\geq 0 $. Our objective is to demonstrate the convexity of the function $ G(z,\alpha)=\sum_{i=0}^{2}\lambda_{i}  F_{2,i}(z)$ over the interval $[0,1]$. Evaluating the second derivative, we get

\begin{equation*}
\dfrac{d^2}{dz^2} G(z,\alpha)=\sum_{i=0}^{2} \lambda_{i}  \dfrac{d^2}{dz^2}  F_{2,i}(z)= 2\alpha \left (\lambda_{2} -2\lambda_{1} +\lambda_{0}   \right )
\end{equation*}

Given this expression and $ \lambda_{2} -2\lambda_{1} +\lambda_{0}\geq 0 $, we can affirm that $ \dfrac{d^2}{dz^2} G(z,\alpha) \geq 0 $.



Presuming the convexity of a set of real values $ \{\beta_{i} \}_{i=0}^{n}$, we take as our induction hypothesis that the expression $ \sum_{i=0}^{n} \beta_{i} F_{n,i}(z) $ is convex.


Now, let's introduce a new set of convex values $ \{\eta_{i}\}_{i=0}^{n+1}$. The goal is to establish the convexity of the function $ \sum_{i=0}^{n+1} \eta_{i} F_{n+1,i}(z) $. To accomplish this, we consider
\begin{equation*} 
G(z,\alpha) =\sum_{i=0}^{n+1} \eta_{i} F_{n+1,i}(z)=\sum_{i=0}^{n+1} \eta_{i}\left [(1-z)F_{n,i}(z)+zF_{n,i-1}(z) \right ],  \end{equation*}
and establish that the function $ \dfrac{d}{dz}G(z,\alpha) $ is  increasing.
\begin{eqnarray*}
\dfrac{d}{dz} G(z,\alpha) &=&\sum_{i=0}^{n+1} \eta_{i} \dfrac{d}{dz} \left [(1-z)F_{n,i}(z)+zF_{n,i-1}(z)  \right ]\\
&=&\sum_{i=0}^{n+1} \eta_{i}\left [-F_{n,i}(z)+(1-z) \dfrac{d}{dz} F_{n,i}(z)+F_{n,i-1}(z)+z \dfrac{d}{dz} F_{n,i-1}(z)  \right ]\\
&=&\sum_{i=0}^{n} ( \eta_{i+1}-\eta_{i})F_{n,i}(z)+(1-z)\sum_{i=0}^{n}\eta_{i} \dfrac{d}{dz} F_{n,i}(z)+z \sum_{i=0}^{n}\eta_{i+1}\dfrac{d}{dz} F_{n,i}(z)
\end{eqnarray*}
Given that $\left \{ \eta_{i}  \right \}_{i=0}^{n+1}$ represents convex data, we observe that $  \eta_{i+1}-2 \eta_{i}+ \eta_{i-1}\geq 0, i=1,\cdots,n-1  $. This inequality implies $  \eta_{i+1}- \eta_{i}\geq \eta_{i}- \eta_{i-1} , i=1,\cdots,n-1  $. In simpler terms, the differences $ \eta_{i}- \eta_{i-1} , i=1,\cdots,n-1 $ form an increasing sequence. By applying Proposition \ref{pro3.2}, we can affirm that the function $ \sum_{i=0}^{n} ( \eta_{i+1}-\eta_{i})F_{n,i}(z) $ is an increasing function.   


Assuming the induction hypothesis and given that $\left \{ \eta_{i} \right \}_{i=0}^{n+1}$ represents convex data, it follows that both $ \sum_{i=0}^{n}\eta_{i}F_{n,i}(z) $ and $ \sum_{i=0}^{n}\eta_{i+1}F_{n,i}(z) $ are convex functions. This implies that their respective derivatives, $ \sum_{i=0}^{n}\eta_{i} \dfrac{d}{dz} F_{n,i}(z)$ and $\sum_{i=0}^{n}\eta_{i+1} \dfrac{d}{dz} F_{n,i}(z) $, are increasing functions.


By the derived results and considering $ z \in [0,1] $, it follows that $ \dfrac{d}{dz} G(z,\alpha) $ is an increasing function. Consequently, this implies $ \dfrac{d^2}{dz^2} G(z,\alpha) \geq 0 $, signifying that $ G(z,\alpha) $ is a convex function.


We observe that if $ q $ is a convex function, then the values $ q\left(\frac{0}{n}\right), q\left(\frac{1}{n}\right), \ldots, q\left(\frac{n}{n}\right) $ form a set of convex data. Consequently, this directly implies that $ R_{n,\alpha}(q;z) $ is also a convex function.
\end{proof}

\end{document}
